\documentclass[12pt,a4paper,oneside]{article}
\usepackage[brazil]{babel}
\usepackage[utf8]{inputenc}
\usepackage{olymp}

\setcounter{problem}{0}

\begin{document}

\begin{problem}{Jogo da memória}{memoria}{2}

O Jogo da Memória é um jogo de tabuleiro clássico focado em memorização e percepção visual.
Um conjunto de cartas com pares idênticos é embaralhado e distribuído em uma grade, todas com a face voltada para baixo.
A cada rodada, o jogador da vez escolhe duas cartas para desvirar:

\begin{itemize}
\vspace{-9pt}
\item Se forem iguais, ele formou um par: elas permanecem viradas para cima e ele ganha um ponto.
\vspace{-9pt}
\item Se forem diferentes: as duas cartas são viradas novamente para baixo após um breve intervalo.
\end{itemize}

\vspace{-9pt}
O jogo termina quando todos os pares forem encontrados e vence quem tiver feito mais pontos.

Considere uma versão de apenas um jogador, em que o próprio jogador continua virando cartas até terminar o jogo (ou cansar de jogar).
Faça um programa para ler suas jogadas e informar quantos pontos ele conseguiu até parar de jogar.

%\Notes
%
%Observações, se tiver.

\InputFile

A entrada começa com dois inteiros, $M$ e $N$, respectivamente o número de linhas e colunas do tabuleiro.
Em seguida virão $M$ linhas, cada uma com $N$ caracteres, representando as cartas do jogo.
Por fim, seguem várias linhas, uma para cada jogada, com quatro números $L_1$, $C_1$, $L_2$, $C_2$ indicando a linha e a coluna das duas cartas que o jogador selecionou na jogada.
As jogadas terminam com uma linha contendo quatro zeros.
Restrições: $1 \leq M, N \leq 100$, $1 \leq L_i \leq M$, $1 \leq C_i \leq N$ e o jogador não vira cartas de pares já formados.

\OutputFile

Seu programa deve escrever o número de pontos obtidos pelo jogador, isto é, o número de jogadas em que as duas cartas escolhidas formaram um par de mesmo valor.

%\Notes

\Example

\begin{example}
\exmp{
2 3
rua
ura
1 1 1 2
2 1 1 2
1 3 2 3
0 0 0 0
}{
2
}%
\end{example}

%Este exemplo corresponde ao da imagem mostrada acima.
Na primeira jogada, o jogador vira as cartas `\texttt{r}' e `\texttt{u}' e então desvira as cartas.\\
Na segunda jogada, ele vira as cartas `\texttt{u}' e `\texttt{u}', formando um par (1 ponto).\\
Na terceira, ele vira as cartas `\texttt{a}' e `\texttt{a}', formando outro par (mais 1 ponto).\\
Então ele para de jogar, ganhando 2 pontos no total.

\begin{example}
\exmp{
4 4
ave.
kkah
the.
trvr
1 1 1 2
2 3 1 1
2 1 2 2
3 1 1 4
0 0 0 0
}{
2
}%
\end{example}

%Comentários sobre o exemplo, se necessário.

\end{problem}

\end{document}